\chapter{Aneurysms}

\section*{Definition and Overview}
An aneurysm is a balloon-like bulge in the wall of an artery, where a weakened area of the vessel stretches outward. Aneurysms can develop in any artery but are most common in the abdominal aorta, in the arteries of the brain (cerebral aneurysms), and behind the knee.

Most aneurysms grow slowly and cause no symptoms, and the greatest danger is that the weakened wall may rupture, causing severe internal bleeding that can be life-threatening. They may also clot, block nearby vessels, or compress surrounding structures. Aneurysms are named by their location and shape, and the type of artery involved determines both the symptoms and the treatment.

\section*{Epidemiology}
Aneurysms affect men more often than women overall, and the risk rises with age. Abdominal aortic aneurysms are most common in men over 65, particularly smokers, and are found in roughly one in twenty older men who are screened.

Cerebral (brain) aneurysms affect roughly 2 to 3\% of the population, and most never rupture; the risk of rupture is higher in women and in people who smoke or have high blood pressure. Popliteal and other peripheral aneurysms are less common and often occur together with aneurysms elsewhere in the body.

\section*{Aetiology and Pathophysiology}
\begin{itemize}
    \item \textbf{Atherosclerosis:} hardening and fatty damage of the artery wall is the commonest cause of aortic and peripheral aneurysms
    \item \textbf{Hypertension:} high blood pressure increases wall stress and speeds up enlargement
    \item \textbf{Smoking:} strongly associated with aneurysm formation and growth
    \item \textbf{Genetics and family history:} a close relative with an aneurysm increases the risk
    \item \textbf{Connective tissue disorders:} conditions such as Marfan syndrome and Ehlers--Danlos syndrome weaken the vessel wall
    \item \textbf{Infection:} rare mycotic aneurysms that develop from infected material carried in the blood
    \item \textbf{Injury:} trauma or previous surgery can damage the vessel wall
\end{itemize}

The wall of a normal artery has elastic layers that withstand pressure. When these layers weaken or are damaged, the wall bulges outward and may progressively enlarge; a larger aneurysm has higher wall tension and therefore a greater risk of rupture.

\section*{Clinical Features}
\begin{itemize}
    \item Most unruptured aneurysms cause no symptoms and are found by chance during imaging
    \item A pulsating feeling or lump in the abdomen (abdominal aortic aneurysm)
    \item Deep, aching pain in the abdomen, back, or flank
    \item A dull pain or pulsatile mass behind the knee (popliteal aneurysm)
    \item Cranial nerve compression from a cerebral aneurysm, causing a drooping eyelid, double vision, or a widened pupil
    \item Sudden severe pain in the chest, abdomen, or back, which suggests expansion or rupture
\end{itemize}

\section*{Differential Diagnosis}
\begin{itemize}
    \item Other causes of abdominal or back pain, including peptic ulcer disease, kidney stones, pancreatitis, musculoskeletal pain, and disc disease
    \item Causes of a sudden severe headache, including migraine, subarachnoid hemorrhage, meningitis, and cervical artery dissection
    \item Chest pain from aortic dissection versus myocardial infarction or pericarditis
    \item A tortuous (curved) aorta or a prominent pulsation in a thin person, which can mimic an aneurysm
    \item Hernias, masses, or tumours that produce a palpable lump
\end{itemize}

\section*{When to Seek Medical Care}
See a doctor promptly if you develop unexplained abdominal, chest, or back pain, especially if you smoke, have high blood pressure, or have a family history of aneurysm. Report a sudden severe headache, double vision, or a drooping eyelid.

\textbf{Call 000 (or local emergency services) for:}
\begin{itemize}
    \item Sudden, severe pain in the abdomen, back, or chest that may indicate rupture
    \item A sudden, extremely severe headache, sometimes described as the worst ever
    \item Fainting, collapse, or loss of consciousness
    \item A rapid heartbeat, clammy skin, or other signs of shock
\end{itemize}

\section*{Diagnosis and Investigations}
\begin{itemize}
    \item \textbf{History and examination:} assessment of risk factors, family history, and palpation of the abdomen or affected limb
    \item \textbf{Ultrasound:} quick and reliable for measuring an abdominal aortic aneurysm, and used for ongoing surveillance
    \item \textbf{CT angiography:} provides detailed images of the aneurysm, its size, and its suitability for repair
    \item \textbf{MRI or MR angiography:} useful for cerebral and complex aneurysms without X-ray exposure
    \item \textbf{Digital subtraction angiography:} used during planning of cerebral aneurysm treatment
    \item \textbf{Screening:} a one-off abdominal ultrasound is offered to men aged 65 and over, particularly smokers
\end{itemize}

\section*{Management}
\begin{itemize}
    \item \textbf{Surveillance:} small aneurysms are monitored with regular ultrasound or CT, and repair is considered when size or growth thresholds are reached
    \item \textbf{Blood pressure control:} medication and lifestyle changes to reduce wall stress
    \item \textbf{Smoking cessation:} essential, because smoking drives enlargement and rupture
    \item \textbf{Lipid-lowering therapy:} statins for underlying atherosclerosis
    \item \textbf{Open surgical repair:} a graft replaces the diseased segment of the vessel
    \item \textbf{Endovascular repair:} a stent graft is placed inside the vessel through a groin artery, avoiding major open surgery
    \item \textbf{Cerebral aneurysm treatment:} endovascular coiling or neurosurgical clipping to exclude the aneurysm from the circulation
\end{itemize}

\section*{Complications}
\begin{itemize}
    \item Rupture, causing life-threatening internal bleeding, with the risk increasing with size
    \item Subarachnoid hemorrhage after cerebral aneurysm rupture, which carries a high risk of brain injury
    \item Vasospasm, stroke, and re-bleeding after subarachnoid hemorrhage
    \item Clot formation within the aneurysm, which can block the vessel or break off and travel (embolism)
    \item Compression of nearby structures by a large aneurysm
    \item Complications of treatment, including bleeding, infection, graft problems, or stroke
\end{itemize}

\section*{Prognosis and Outcomes}
Small aneurysms often remain stable for years and may never rupture. When repaired electively at an appropriate size, the outlook is generally good and the risk of surgery is relatively low. Ruptured aneurysms are a medical emergency, and many people do not survive despite urgent treatment. After successful repair, most people return to normal activities, although some need long-term blood pressure control and surveillance of the rest of their blood vessels.

\section*{Prevention and Self-Care}
\begin{itemize}
    \item Do not smoke, as smoking is the strongest modifiable risk factor
    \item Keep blood pressure and cholesterol in a healthy range
    \item Maintain a healthy weight and be physically active
    \item Discuss screening with a doctor if you are a man aged 65 or over, or if an aneurysm runs in your family
    \item If you have a connective tissue disorder such as Marfan syndrome, follow specialist surveillance advice
    \item Attend scheduled imaging follow-ups so that any growth is detected early
    \item Seek medical advice promptly for new or unusual pain in the chest, abdomen, or back
\end{itemize}

\section*{Sources and Further Reading}
The following resources provide reliable, peer-reviewed and patient-orientated information on aneurysms.

\begin{itemize}
    \item NHS. \textit{Overview of aneurysms, screening, and treatment.} https://www.nhs.uk/conditions/
    \item Mayo Clinic. \textit{Guide to the types, diagnosis, and treatment of aneurysms.} https://www.mayoclinic.org/
    \item Centers for Disease Control and Prevention. \textit{Information on aortic aneurysms and risk factors.} https://www.cdc.gov/
    \item MSD Manual. \textit{Professional and patient information on aneurysms and their management.} https://www.msdmanuals.com/
    \item MedlinePlus. \textit{Health information on aneurysms.} https://medlineplus.gov/
    \item NICE. \textit{Clinical guidance on the diagnosis and treatment of abdominal aortic aneurysms.} https://www.nice.org.uk/
\end{itemize}
